\documentclass[11pt]{article}
\usepackage[T1]{fontenc}
\usepackage[utf8]{inputenc}
\usepackage{lmodern,microtype}
\usepackage[margin=1in]{geometry}
\usepackage{amsmath,amssymb,amsthm,mathtools}
\usepackage{booktabs,array,longtable}
\usepackage{xcolor}
\usepackage{hyperref}
\usepackage{enumitem}
\hypersetup{colorlinks=true,linkcolor=blue!45!black,citecolor=blue!45!black,
 urlcolor=blue!45!black,pdftitle={Centered Pell gaps and adaptive congruence moduli},
 pdfauthor={Research contribution prepared for Vladimir Reshetnikov}}
\setlength{\emergencystretch}{2em}
\setlist{itemsep=3pt,topsep=5pt}
\newcommand{\Z}{\mathbb Z}
\newcommand{\N}{\mathbb N}
\newcommand{\Zp}{\mathbb Z_{>0}}
\newcommand{\Con}{\operatorname{Con}}
\newcommand{\pc}{\operatorname{pc}}
\newcommand{\ini}{\operatorname{in}}
\newtheorem{theorem}{Theorem}[section]
\newtheorem{lemma}[theorem]{Lemma}
\newtheorem{proposition}[theorem]{Proposition}
\theoremstyle{remark}
\newtheorem{remark}[theorem]{Remark}
\title{Centered Pell gaps and adaptive congruence moduli\\
\large Explicit improvements for Diophantine consistency equations}
\author{Research contribution prepared for Vladimir Reshetnikov\\
\small ProveIt project}
\date{23 October 2026}
\begin{document}
\maketitle
\begin{abstract}
We give two concrete improvements to the latest supplied ProveIt constructions.
A fixed polynomial in thirteen strictly positive integer unknowns has exact
degree $30$ and a $93$-operation arithmetic circuit, improving the previous
degree $34$ at the same circuit cost. A second fixed polynomial in twelve
strictly positive integer unknowns has exact degree $110$ and $100$ operations,
improving the preceding degree-$118$, $104$-operation construction on both
measures. Their lack of positive integer solutions is equivalent to
$\Con(\mathrm{ZFC})$ by ordinary mathematical proofs formalizable in ZFC,
composed with the pinned native compiler and refutation calibration.

The first improvement uses an already recovered divisibility condition to
replace a fifth-power congruence modulus by a lower-degree adaptive modulus.
The possible common divisor three is handled exactly. The second uses the
small positive gap between two auxiliary Pell coordinates to merge the
strong Pell root and auxiliary witness, and then centers the strong root
before taking the norm. We prove that the exceptional zero denominator is
impossible. The archive contains the four actual coefficient numerals, their
reconstruction recipe, literal polynomial circuits, independent exact checks,
and full proofs of the new reductions. These results do not lower the existing
nine-positive-variable frontier or establish global minimality. They are
specific improvements, not a claimed completely new consistency encoding.
\end{abstract}
\setcounter{tocdepth}{1}
\tableofcontents
\newpage

\section{Outcome, comparison, and scope}

\subsection{What is improved}

The comparison snapshot is ProveIt commit
\texttt{a2cd79b9120fc75f48dcbef6e38d0f8b18d64d49}, which places sources
60--67 beside the earlier packages. In particular, source 61 already gives
the strong twelve- and thirteen-positive-variable points used as our direct
baselines \cite{source61}. The older collection summary alone would give a
misleadingly weak baseline.

\begin{center}\small
\begin{tabular}{@{}lrrrr@{}}
\toprule
Construction & Unknowns & Degree & Operations & Residuals\\
\midrule
Source 61 & 13 & 34 & 93 & 4\\
Adaptive modulus, this report & 13 & \textbf{30} & \textbf{93} & 4\\
Source 61, product form & 12 & 118 & 104 & 3\\
Centered gap, this report & 12 & \textbf{110} & \textbf{100} & 3\\
Source 61, sum of squares & 12 & 172 & 103 & 3\\
Centered gap, sum of squares & 12 & \textbf{164} & \textbf{99} & 3\\
\bottomrule
\end{tabular}
\end{center}

The first degree reduction is $4/34\approx11.76\%$. The primary twelve-variable
degree reduction is $8/118\approx6.78\%$, accompanied by four fewer binary
operations. All operations count: addition, subtraction, multiplication,
squaring, and multiplication by a fixed numeral each cost one gate. Fixed
integer leaves and reuse of a previously computed register are free. This is
an arithmetic-circuit measure, not bit complexity.

The current repository also contains nine-positive-variable constructions,
ten-integer constructions, and much lower-degree equations with more
variables. None is displaced merely by the table above. The best results of
this iteration are at the twelve- and thirteen-positive-variable slices.
In particular, source 64's $41$-integer quartic and source 62's $35$-integer
sextic and $30$-integer octic remain better low-degree integer choices
\cite{source62,source64}. A proposed combination of those sources was checked
and found to double-count an already eliminated register; no additional
variable reduction is claimed.
Source 62 also retains a cheaper thirteen-positive-variable point of degree
$114$ in $83$ operations. Our thirteen-variable result improves degree at
the fixed $93$-operation source-61 baseline; it does not dominate every
thirteen-variable tradeoff.

\subsection{What ``provably equivalent'' means here}

Fix the usual arithmetic proof predicate for finite ZFC refutations used by
the repository, and write
\[
 \Con(\mathrm{ZFC})\equiv
 \neg\exists p\,\operatorname{Prf}_{\mathrm{ZFC}}(p,\bot).
\]
The intended theorem for our primary polynomial $P_{12}$ is
\begin{equation}
 \mathrm{ZFC}\vdash
 \left(\neg\exists\mathbf x\in\Zp^{12}\;P_{12}(\mathbf x)=0\right)
 \longleftrightarrow \Con(\mathrm{ZFC}).
 \label{eq:target}
\end{equation}
We establish the new reductions in ordinary mathematics, with explicit
witness recovery in both directions. The native coefficient/compiler theorem
and the fixed K215/refutation calibration are inherited mathematical inputs,
identified in Section~\ref{sec:dependencies}; no new end-to-end Lean, Rocq,
or object-language ZFC certificate is supplied. Exact symbolic checks verify
the finite algebra and literal circuits. They do not replace the unbounded
compiler proof. Equation~\eqref{eq:target} asserts an equivalence, not either
consistency or inconsistency.

The arithmetic construction uses the repository's native kernel. It does
not use Jones's universal polynomial as its finalizer. Alternative source
architectures were considered, but no complete replacement with a better
measured endpoint was obtained. No claim of literature priority is made for
elementary norm or Pell identities; the new contribution is the precise
combination and its verified effect on these fixed equations.

\section{Four fixed numerals, with no hidden parameter}
\label{sec:coefficients}

Every equation here uses the same complete source-44 compact tuple
\cite{source44}:
\[
 (\beta,K_0,M_C,M_F),\qquad \beta=2^{590069}-1.
\]
The archive contains their actual bytes in four gzip files. A decompressed
byte string $b_0,\ldots,b_{m-1}$ denotes the integer
$\sum_{j<m}b_j256^j$. The port named $K_0$ is the payload named
\texttt{K}; the port $M_F$ is \texttt{MF\_source}, which already includes
the shift by $\beta$. One must use all four together.

\begin{center}
\begin{tabular}{@{}lrrl@{}}
\toprule
Port & Bit length & Raw bytes & Payload\\
\midrule
$\beta$ & 590069 & 73759 & \texttt{beta.bin.gz}\\
$K_0$ & 592884 & 74111 & \texttt{K.bin.gz}\\
$M_C$ & 590069 & 73759 & \texttt{MC.bin.gz}\\
$M_F$ & 590070 & 73759 & \texttt{MF\_source.bin.gz}\\
\midrule
Total & & 295388 &\\
\bottomrule
\end{tabular}
\end{center}

These are fixed positive integers, not extra quantified variables and not
constants chosen according to the truth of consistency. The supplied
\texttt{upstream/build\_coefficients.py} reconstructs them from the supplied
430-row transition table, reconstructs the 943-tile language, and checks the
finite packing certificate. Its generated raw numeral hashes are compared
against the delivered binding by \texttt{verify.py}. This recipe is a
finite computation, not an unbounded proof or halting search.

The numerical magnitude is still large. This iteration reduces arithmetic
degree and circuit size, not coefficient height. Naming a large numeral is
legitimate in a polynomial but should not be advertised as reducing its
binary description length. The supplied payloads and recipe make that cost
visible and reproducible.

\section{Common polynomial registers}
\label{sec:common}

For the thirteen-variable construction supply precisely
\[
 j,F,\alpha,f,i,T,X,Y,r,\eta,\zeta,u,\gamma\in\Zp.
\]
For the twelve-variable construction replace the pair $(f,u)$ by one
positive input $n$:
\[
 j,F,\alpha,i,T,X,Y,r,\eta,\zeta,n,\gamma\in\Zp.
\]
All subsequent letters are polynomial abbreviations. They add no unknowns.
For the thirteen-variable equation put $J=3j+1$; for the twelve-variable
equation put $J=6j+1$. Define
\begin{align}
 q&=\beta J+1,& C&=q-F-\alpha,& E&=XY,\\
 k&=\eta+\zeta,& c&=kY+\eta,& a&=E+Y,\\
 H&=4a+3,& \Delta&=a^2+H=(a+2)^2-1,& I&=ic^2,\\
 \ell&=EkY,& b_0&=X+ac,\\
 R&=\bigl((q-1)(q-F)+\alpha\bigr)(q^2-1)
                  +(M_C+qM_F)J.&&&\label{eq:R}
\end{align}
The index defect, adaptive modulus, and its numerator are
\begin{align}
 r_0&=k-R-1,& L&=qY,& U&=qL=q^2Y,\\
 D_*&=Ub_0+H(U\gamma+LX-r_0).&&&\label{eq:Dstar}
\end{align}
For the first carrier put
\begin{align}
 W&=(q-1)k+1,& M_0&=q^3(q-1),\\
 B_t&=(qK_0+X)C-qF,\\
 Z&=M_0(\ell+k+r)+(q-1)Y-q^3WB_t.&&&\label{eq:Z}
\end{align}
Here and below $qF$ is multiplication by the supplied input $F$.
The two carrier residuals are
\begin{equation}
 A_0=Z^2-M_0^2\bigl(\ell(\ell+k)+1\bigr),\qquad
 A_1=D_*^2-U^2(1+\Delta c^2).
 \label{eq:carriers}
\end{equation}
Every power in these definitions is a fixed nonnegative integer power.
There are no variable exponents, divisions, roots, digit functions, or
implicit witnesses in the actual polynomials.

\section{The thirteen-variable equation}
\label{sec:thirteen}

With the common registers and $J=3j+1$, define
\begin{align}
 Q_f&=\Delta(f^2-1),& V_f&=f^2(cT-R)-c,\\
 A_s&=f^2-\Delta I^2-1,\\
 A_a&=(V_f+u)^2-Q_fu(2V_f+u)-1.
\end{align}
Our first equation is
\begin{equation}
 \boxed{\;P_{13}=A_0^2+A_1^2+A_a^2+A_s^2=0,\quad
                 \mathbf x\in\Zp^{13}.\;}
 \label{eq:P13}
\end{equation}

\begin{theorem}
Equation~\eqref{eq:P13} has a positive integer solution exactly when the
pinned native consistency kernel has one. Consequently its unsolvability
is ZFC-provably equivalent to $\Con(\mathrm{ZFC})$. Its exact total degree
is $30$, and the supplied literal evaluator has $93$ live binary operations.
\end{theorem}

We give the new argument in detail because a careless cancellation in its
congruence step would invalidate the improvement.

\subsection{Recover the first scale and the sign of the transport}

At a positive zero all four residuals vanish. Since $M_0>0$, the equation
$A_0=0$ says $Z/M_0$ is a rational square root of the integer
$\ell(\ell+k)+1$. A rational square root of an integer is an integer.
Thus $M_0\mid Z$. Reducing this divisibility modulo $q^3$, and then
modulo $q-1$, gives
\[
 q^3\mid Y,\qquad q-1\mid B_t.
\]
The reductions use $\gcd(q,q-1)=1$ and $W\equiv1\pmod{q-1}$.
Set $s=Y/q^3>0$, $t=B_t/(q-1)$, and $\tau=Z/M_0$. Then
\begin{equation}
 \tau^2=\ell(\ell+k)+1,\qquad
 \tau=\ell+k+r+s-Wt.\label{eq:tau}
\end{equation}
Because $k\ge2$, one has $|\tau|<\ell+k$. Hence $Wt>r+s>0$,
so $t>0$ and $B_t>0$. But $C\le0$ would give
$(qK_0+X)C-qF<0$. We have proved $C>0$ before using any divisibility
of $X$. The inherited compact-mask bounds now give
\begin{equation}
 q\ge16,\qquad 3q<R<q^4-2.\label{eq:outerbounds}
\end{equation}

Since $590069$ is odd, $\beta\equiv1\pmod3$. Our restriction on $J$
gives $q\equiv2\pmod3$. Since $q\mid Y$,
$H=4Y(X+1)+3\equiv3\pmod q$, and therefore
\begin{equation}
 \gcd(q,H)=1.\label{eq:qH}
\end{equation}

\subsection{Recover the main Pell index before the second scale}

The numerator $D_*$ is positive: $r_0<k$, $c>kY$, and $U\ge q^5$ imply
\[
 D_*>Uac-Hk>k(UaY-H)>0.
\]
The last inequality follows immediately from $q\ge16$, $Y\ge q^3$,
and $H=4a+3$. Applying rational-square integrality to $A_1=0$ gives
an integer $g_0=D_*/U>0$ with
\begin{equation}
 g_0^2-\Delta c^2=1,\qquad U\mid H(LX-r_0).
 \label{eq:adaptivecongruence}
\end{equation}

Write $A=a+2$ and $P_0=2XY^2+1$. Use the usual Pell sequences
\[
 \chi_b(m)+\psi_b(m)\sqrt{b^2-1}
       =(b+\sqrt{b^2-1})^m.
\]
The source Pell classification supplies positive indices $p,n_0$ with
\[
 c=\psi_A(p),\quad g_0=\chi_A(p),\quad
 k=2\psi_{P_0}(n_0),\quad |\tau|=\chi_{P_0}(n_0).
\]
Here $c>1$ excludes $p=1$, because $\psi_A(1)=1$.
The case $n_0=1$ would give $k=2$, $\eta=\zeta=1$, and $c=2Y+1$,
whereas $p\ge2$ gives $c\ge2A\ge4Y+4$. Therefore $n_0\ge2$, and
\begin{equation}
 c>kY>8XY^3\ge8q^9>2R,\qquad n_0<p,\qquad c>2p.
 \label{eq:earlybounds}
\end{equation}
The comparison $n_0<p$ uses $P_0>A$ and $c>k$.

The normalized strong equation $A_s=0$, with $I=ic^2$, gives
$f=\chi_A(m)$, $I=\psi_A(m)$, $c\mid m$, and $m>2p$.
For clarity, strong divisibility first gives $p\mid m$; writing $m=pv$,
the composition formula yields $I/c=\psi_{g_0}(v)$. Since $c\mid I/c$
and $g_0^2\equiv1\pmod c$, its recurrence gives
$\psi_{g_0}(v)\equiv v g_0^{v-1}\pmod c$.
The residue $g_0$ is invertible modulo $c$, so $c\mid v$.
Thus $m\ge pc>2p$.

Now $Q_f=(\Delta I)^2$ and $A_a=0$ yield
\[
 (\Delta I)^2V_f^2-\bigl((\Delta I)^2-1\bigr)|V_f+u|^2=1.
\]
Furthermore $V_f\equiv-c\pmod f$ and $V_f\equiv-R\pmod c$.
Together with~\eqref{eq:earlybounds}, these are precisely the hypotheses
of the retained signed auxiliary-index lemma \cite{source61}. It gives
$p=R$. A complete retained proof is supplied in
\texttt{upstream/SIGNED\_INDEX\_LEMMA.md}. This argument has not assumed
$q\mid X$ or reconstructed the missing index quotient.

\subsection{The adaptive modulus and its common divisor three}

The old source uses $U=q^5$ and $L=q^4$. We use $U=q^2Y$ and $L=qY$.
This decreases their formal degrees from $5,4$ to $3,2$ while preserving
the size needed for index recovery. The precise reason is the already
proved divisibility $q^3\mid Y$.
Adaptive and successive congruence recovery already appear in sources 64
and 65 \cite{source64,source65}. The contribution here is their exact
integration with the source-61 thirteen-variable carrier, using the modulus
$q^2Y$ and the following cofactor argument.

Put
\[
 h_3=\gcd(H,U)=\gcd(H,Y)=\gcd(3,Y)\in\{1,3\},
 \qquad \mu=\frac{qY}{h_3}.
\]
Since $\gcd(q,h_3)=1$, dividing $Y$ by $h_3$ preserves $q^3\mid Y/h_3$.
In particular
\begin{equation}
 \mu\ge q^4>R.\label{eq:mu}
\end{equation}
Canceling exactly $h_3$ in~\eqref{eq:adaptivecongruence} shows
$U/h_3\mid LX-r_0$, and reduction modulo $\mu=L/h_3$ gives
$\mu\mid r_0$. The first Pell recurrence gives
\[
 k\equiv2n_0\pmod{4XY^2},
\]
and $\mu\mid4XY^2$. Thus $\mu\mid2n_0-R-1$.
Since $1\le n_0<p=R<\mu$, the absolute value is smaller than $\mu$;
therefore
\begin{equation}
 2n_0=R+1.\label{eq:exactfirstindex}
\end{equation}

It follows that $r_0=k-2n_0$ is a positive multiple of $4XY^2$.
Because $q^2\mid Y$, this is also a multiple of $U=q^2Y$.
Equation~\eqref{eq:adaptivecongruence} now gives $U\mid HLX$;
as $U=qL$, this implies $q\mid HX$. By~\eqref{eq:qH}, $q\mid X$.
The important distinction is that the common factor three is removed from
an already divisible scale. Merely estimating $qY/3\ge q^4/3$ would be
too weak and would miss the actual argument.

\subsection{Restore all deleted positive witnesses}

Define
\[
 w=X/q,\qquad h=r_0/E,\qquad
 \gamma_{\mathrm{old}}=\gamma+w-r_0/U.
\]
These are integers; the first two are positive. The main root satisfies
$g_0=X+ac+H\gamma_{\mathrm{old}}$. Since
$g_0>(a+1)c$ and $c>X$, it follows that
$H\gamma_{\mathrm{old}}>c-X>0$, so the third is positive as well.

Now $B_t=q[(K_0+w)C-F]$. Its positivity and its divisibility by $q-1$
give a positive integer $z_{\mathrm{new}}$ with
\[
 (K_0+w)C-F=(q-1)z_{\mathrm{new}}.
\]
Set $z_{\mathrm{old}}=z_{\mathrm{new}}+1$; then
\[
 (K_0+w)C+(q-F)-z_{\mathrm{old}}(q-1)=1.
\]
With $T_{\mathrm{old}}=fT$, the roots $|\tau|,|V_f+u|$, and the
restored index quotient $h$, this is the full retained normalized source
interface. To use the ordinary strong-equation interface instead, put
$i_{\mathrm{old}}=\Delta i$; then
$(i_{\mathrm{old}}c^2)^2=\Delta(f^2-1)$.
Source soundness gives a finite ZFC refutation.

\subsection{Construct positive new witnesses from a refutation}

A finite refutation supplies the canonical compact witness of the inherited
compiler. It has $q=B^N$, $B=2^{590069}$, and
$J=1+B+\cdots+B^{N-1}$ for odd $N>1$. Hence $J\equiv1\pmod3$,
$j=(J-1)/3>0$, and $J$ is odd. Its other properties include
\[
 X=2^R,\quad w=X/q>q,\quad Y=sq^3,\quad R\equiv3\pmod4,
 \quad \gamma_{\mathrm{old}}>w.
\]
The retained gap estimates and auxiliary construction are documented in
source 58 and source 61 \cite{source58,source61}. In particular $r_0$ is
a positive multiple of $4XY^2$, so
\[
 \gamma=\gamma_{\mathrm{old}}-w+r_0/U>0
\]
is integral and makes $D_*=Ug_0$. For the first carrier choose
\[
 t=q(z_{\mathrm{old}}-1),\qquad
 r=Wt+\tau-\ell-k-s.
\]
Here $z_{\mathrm{old}}\ge2$ canonically, $t\ge q$, and
$\tau>\ell$, $k>s$, so
$r>qW-k-s>0$. Direct substitution gives $Z=M_0\tau$.
Finally reselect the canonical squared-modulus auxiliary pair, retaining
positive $f,i,T,u$. This gives $A_s=A_a=0$. All supplied coordinates are
positive, completing the converse. No transformation of an arbitrary old
witness is claimed; canonical witness reselection is sufficient for the
existence equivalence.

\section{A centered two-root elimination}
\label{sec:norm}

The following local lemma is the additional ingredient for twelve variables.
It applies before any compiler reasoning and is useful independently.
The two-root norm identity and the strategy of centering a root are already
used in source 62, Sections 2--3 of its positive-thirteen proof
\cite{source62}. That construction centers the first Pell root at $\ell$.
Here we merge the strong root with the auxiliary gap and center it at
$(a+2)I$. The distinct center, its positivity bound, and the exceptional-case
exclusion produce the twelve-variable improvement below.

\begin{lemma}[Centered Pell-gap norm]
Let $a\ge2$, $c\ge3$, and $i\ge1$ be integers, and let $b$ be any integer.
Put
\begin{align}
 \Delta&=(a+2)^2-1,& I&=ic^2,&D&=1+\Delta I^2,\\
 Q&=\Delta^2I^2,&d&=Q-1,&f_0&=(a+2)I,\\
 e&=1-I^2,&V&=Db-c.
\end{align}
For an integer $n$ define
\begin{align}
 x&=n+f_0+V,& E_a&=QV^2-1,\\
 N&=\bigl(d(x^2-D)-E_a\bigr)^2-4dDE_a.
 \label{eq:Nliteral}
\end{align}
If $N=0$, there exist positive integers $f,u$ with
\begin{equation}
 f^2=D,\qquad QV^2-d(V+u)^2=1.
 \label{eq:restorednorms}
\end{equation}
Conversely, if a positive solution of~\eqref{eq:restorednorms} satisfies
$u>f_0$, then $n=f-f_0+u$ is positive and makes $N=0$.
\end{lemma}

\begin{proof}
The centered error is exactly $D-f_0^2=e=1-I^2$. Define the proof registers
\[
 B_n=2dx,\qquad H_n=d(n^2+2Vn+e)-V^2+1.
\]
Two elementary identities give
\begin{align}
 N&=\bigl(d(x^2+D)-E_a\bigr)^2-B_n^2D,\\
 d(x^2+D)-E_a&=H_n+B_nf_0.
\end{align}
Consequently
\begin{equation}
 N=(H_n+B_nf_0)^2-B_n^2D
   =H_n^2+2B_nf_0H_n-B_n^2e.
 \label{eq:Ncentered}
\end{equation}

We must exclude $B_n=0$, rather than divide without justification.
Since $d>0$, that case would give $x=0$ and
\begin{align*}
 H_n&=dD-QV^2+1\\
    &=Q(D-V^2)-D+1\\
    &=\Delta I^2\bigl(\Delta(D-V^2)-1\bigr)\ne0.
\end{align*}
The last inequality uses $\Delta>1$: no integer $D-V^2$ can have
$\Delta(D-V^2)=1$. Hence $N=H_n^2\ne0$, a contradiction.

For $N=0$, therefore,
$f_*=(H_n+B_nf_0)/B_n$ is a rational square root of the integer $D$.
It is an integer. Put $y_*=x-f_*$. Expanding the preceding identities gives
\[
 QV^2-dy_*^2=1.
\]
Also $|V|>1$. Indeed $D>c+1$; for $b\le0$, $V\le-c$, and for
$b\ge1$, $V\ge D-c>1$. Therefore
\[
 d(y_*^2-V^2)=V^2-1>0,
\]
so $|y_*|>|V|$. The positive integers
$f=|f_*|$ and $u=|y_*|-V$ satisfy~\eqref{eq:restorednorms}.

Conversely put $n=f-f_0+u$. Then $x=f+(V+u)$, and substituting these
two roots into the uncentered norm gives $N=0$. The stated gap bound
ensures $n>0$.
\end{proof}

\begin{remark}
The soundness direction does not require $n>0$ and allows both signs of
the intermediate roots. Positivity is recovered afterwards. The converse
does require room for the subtraction of $f_0$; this is supplied by the
following canonical Pell estimate, not by an unchecked change of variables.
\end{remark}

\begin{lemma}[Canonical gap]
For $M\ge2$ and odd $R\ge5$,
\[
 \psi_M(R)-\frac{\chi_M(R)}M
   =\frac{\psi_M(R-1)}M\ge8M^2-4.
\]
Both terms on the left are integers.
\end{lemma}
\begin{proof}
The recurrence identity $\chi_M(R)=M\psi_M(R)-\psi_M(R-1)$ proves the
equality. For odd $R$, $\chi_M(R)$ is a multiple of $M$ by the odd-index
Chebyshev recurrence. The sequence $\psi_M(t)$ increases for $t\ge0$;
$R-1\ge4$ and $\psi_M(4)=8M^3-4M$ give the bound.
\end{proof}

The proof uses a very large canonical auxiliary gap to pay for the small
polynomial center. It does not estimate the size of a ZFC refutation or
assume that one exists.

\section{The twelve-variable equation}
\label{sec:twelve}

Use the twelve positive inputs from Section~\ref{sec:common}, now with
$J=6j+1$. Retain $A_0,A_1$ from~\eqref{eq:carriers}. Put
\begin{align}
 D&=1+\Delta I^2,& Q&=\Delta(D-1)=\Delta^2I^2,&d&=Q-1,\\
 f_0&=(a+2)I,&V&=D(cT-R)-c,&x&=n+f_0+V,\\
 E_a&=QV^2-1,&N&=\bigl(d(x^2-D)-E_a\bigr)^2-4dDE_a.
\end{align}
The primary equation is
\begin{equation}
 \boxed{\;P_{12}=(1+A_0)(1+A_1)(1+N)-1=0,
                  \quad\mathbf x\in\Zp^{12}.\;}
 \label{eq:P12}
\end{equation}
This is a single fixed integer polynomial, fully defined by the displayed
registers and the four actual numerals in Section~\ref{sec:coefficients}.

\begin{theorem}
Equation~\eqref{eq:P12} has a positive integer zero if and only if there is
a finite ZFC refutation. The proof is formalizable in ZFC using the retained
calibration. The equation has exact total degree $110$ and a $100$-operation
literal circuit.
\end{theorem}

\subsection{The product finalizer is exact}

Because $J$ and $\beta$ are odd, $q$ is even. Thus $M_0$ and $U$ are
even, and
\[
 A_0\equiv Z^2\pmod4,\qquad A_1\equiv D_*^2\pmod4.
\]
Also $N\equiv\bigl(d(x^2-D)-E_a\bigr)^2\pmod4$ directly from its
definition. If $P_{12}=0$, its three integer factors multiply to $1$,
so each is $+1$ or $-1$. A factor $-1$ would require its residual to be
$-2\equiv2\pmod4$, which is not a square. All factors are therefore
$+1$:
\begin{equation}
 P_{12}=0\quad\Longleftrightarrow\quad A_0=A_1=N=0.
 \label{eq:productexact}
\end{equation}
No nonnegativity of the residuals is assumed. For a more familiar but
higher-degree finalizer, the delivered companion is
$A_0^2+A_1^2+N^2=0$.

\subsection{Soundness and completeness}

Given a positive zero of $P_{12}$, the local norm lemma with $b=cT-R$
restores positive $f,u$ satisfying the strong and auxiliary residuals.
The outer registers are unchanged. Set $j_{13}=2j$; then
$3j_{13}+1=6j+1$, so the recovered tuple is a zero of the proved
thirteen-variable polynomial. Its soundness gives a finite ZFC refutation.

Conversely take the canonical source tuple for a refutation. Its repunit
$J$ is odd and $1$ modulo $3$, hence $1$ modulo $6$; padding gives $J>1$.
Thus $(J-1)/6$ is a positive integer. Use the adaptive-carrier construction
above. For the normalized strong witness set $M=\Delta I$. The canonical
auxiliary pair is
\[
 V=\chi_M(R)/M,\qquad y=\psi_M(R),
\]
with $R\equiv3\pmod4$ and $R>3q\ge48$. The squared-modulus source
construction supplies positive integer $T$ with $V=D(cT-R)-c$.
The canonical gap lemma gives
\[
 u=y-V\ge8M^2-4>f_0,
\]
because $f_0=(a+2)I<\Delta I=M$. Hence
\[
 n=f-f_0+u>0
\]
is an admissible new input. The norm lemma gives $N=0$, and the first two
residuals already vanish. This proves the converse and therefore
Equation~\eqref{eq:target}.

\section{Exact degrees and circuit accounting}
\label{sec:degree}

All supplied unknowns have degree one and the four fixed numerals degree
zero. The actual degree is often lower than naive degree propagation
through the evaluation graph. The verifiers check the cancellations rather
than reporting such propagation as an exact result.

Put $q_*=3\beta j$ or $6\beta j$ as appropriate, and write
\[
 k_*=\eta+\zeta,\quad a_*=XY,\quad c_*=k_*Y,
 \quad\Delta_*=X^2Y^2,\quad I_*=ik_*^2Y^2,
\]
\[
 R_*=q_*^3(q_*-F),\qquad
 B_{t,*}=(K_0q_*+X)(q_*-F-\alpha)-q_*F.
\]
The carrier leading forms, checked by full exact expansion of the
low-degree residuals, are
\begin{align}
 \ini_{15}(A_0)&=-2q_*^8k_*^2XY^2B_{t,*},\label{eq:A0lead}\\
 \ini_{13}(A_1)&=8q_*^2k_*X^2Y^4(q_*^2Y\gamma+R_*).
 \label{eq:A1lead}
\end{align}
In particular the sign before $R_*$ in~\eqref{eq:A1lead} is positive.
The other thirteen-variable residual degrees are $13$ and $14$.
Only $A_0^2$ reaches degree $30$, giving
\[
 \ini_{30}(P_{13})=4q_*^{16}k_*^4X^2Y^4B_{t,*}^2.
\]
The coefficient of $j^{16}\eta^4X^4Y^4\alpha^2$ is
$4(3\beta)^{16}>0$. Specializing the actual fixed numerals therefore
cannot cancel this top degree.

For the centered norm, the convenient degree certificate is
Equation~\eqref{eq:Ncentered}, not the shorter literal evaluator.
The degrees are
\begin{center}
\begin{tabular}{@{}lrrrrrrrr@{}}
\toprule
Register & $D$ & $Q,d$ & $V$ & $f_0$ & $e$ & $B_n$ & $H_n$ & $N$\\
\midrule
Degree & 14 & 18 & 18 & 7 & 10 & 36 & 37 & 82\\
\bottomrule
\end{tabular}
\end{center}
The three terms $H_n^2$, $2B_nf_0H_n$, and $-B_n^2e$ have degrees at
most $74,80,82$. The last is uniquely highest, and
\begin{equation}
 \ini_{82}(N)=4\Delta_*^6I_*^{10}R_*^2\ne0.
 \label{eq:Nlead}
\end{equation}
Thus the product finalizer has exact degree $15+13+82=110$, and the
sum-of-squares companion has degree $164$. A specific nonzero coefficient
in $P_{12}$ is
\begin{equation}
 [j^{18}F^2\alpha\gamma i^{10}X^{16}Y^{39}\eta^{23}]P_{12}
       =64(6\beta)^{18}>0.
 \label{eq:degreewitness}
\end{equation}
Its total exponent is $110$. The analogous top coefficient in the
sum-of-squares companion is
\[
 [j^{12}F^4i^{20}X^{24}Y^{64}\eta^{40}]P_{12,\mathrm{sos}}
       =16(6\beta)^{12}>0.
\]

The literal primary circuit evaluates $N$ by the shorter
Equation~\eqref{eq:Nliteral}. The proof certificate uses the algebraically
equal centered expression to expose degree cancellation. Neither expression
introduces an additional witness. The graph is the exact polynomial
definition; its registers are ordered abbreviations, and every gate and
input is live.

\section{Domains: positive, nonnegative, and unrestricted}

The small variable counts above are over $\Zp$. They must not be relabeled
as unrestricted integer counts. For nonnegative integers simply replace
each of the twelve inputs by $x+1$. The delivered
\texttt{circuits/natural12.json} has degree $110$ and $112$ gates.

For an entirely unrestricted integer equation, define
\[
 \theta(a,b,c)=1+a^2+b^2+c(c+1).
\]
This takes exactly the positive integer values. Nonnegativity is immediate.
For surjectivity, Legendre's three-square theorem applied to $4m+1$
gives one odd and two even squares; write
\[
 4m+1=(2a)^2+(2b)^2+(2c+1)^2,
\]
so $m=a^2+b^2+c(c+1)$. This familiar conversion is also recorded in
Lemma 6.1 of Bayer et al.~\cite{bayer}.

Applying $\theta$ to each input gives a fully specified equation in $36$
unrestricted integers, of exact degree $220$ and $184$ gates. Its
unsolvability has the same ZFC equivalence. It is included to remove any
domain ambiguity, not as a new integer frontier point. Its degree is exact:
the homogeneous leading substitution maps each original variable to a sum
of three distinct squares; setting the second and third members of each
triple to zero reduces to the injective substitution $x_i\mapsto a_i^2$.

\section{A tempting shortcut that fails}
\label{sec:barrier}

One might try to replace $X=wq$, $Y=sq^3$ by $X=wq^2$, $Y=sq^2$,
preserving $XY\ge q^4$ while lowering the first carrier's degree. The
compiler does not depend only on the product $XY$.

For $q=2^t$, its half-binomial encoding has a packed index
\[
 R=(q^2-S)(q^2-1)+T,\qquad \pc(T)=t+2.
\]
The relevant inverse-population lemma gives $\pc(R)\le3t+2$,
with equality exactly when the required masks vanish. The corresponding
ordinate valuation is $\pc(R)-2$. Thus $q^3\mid Y$ enforces the full
threshold, whereas $q^2\mid Y$ does not. This role of the cubic scale is
explicit in the retained source notes, cited with exact blobs in
\texttt{research/SCALE\_BARRIER.md}.

An exact local negative control is
\[
 q=16,\quad t=4,\quad S=32,\quad T=175,\quad R=57295.
\]
Here $\pc(T)=6=t+2$, $R\equiv3\pmod4$, but $\pc(R)=13<14=3t+2$
and $S\mathbin{\&}T=32\ne0$. The weakened threshold $2t+2=10$ accepts
this invalid local mask configuration. This is a counterexample to the
weakened packing predicate, not a purported zero of the full consistency
polynomial. Moving a factor of $q$ to $X$ cannot repair this missing
population condition by a size estimate alone.

A substantial further reduction will need a different local packing
predicate, a lower-cost way to enforce that exact threshold, or a source
formalism whose complete arithmetic compilation is smaller. Merely citing
a smaller universal machine or the DPRM theorem does not supply the missing
fixed polynomial or its coefficient binding.

\section{Dependencies, verification, and integration}
\label{sec:dependencies}

\subsection{The proof boundary}

The logical chain is
\[
 \begin{gathered}
 \exists\mathbf x\in\Zp^{12}:P_{12}=0\\
 \Longleftrightarrow\ \text{the normalized native kernel has a positive witness}\\
 \Longleftrightarrow\ \exists p\,\operatorname{Prf}_{\mathrm{ZFC}}(p,\bot).
 \end{gathered}
\]
The first equivalence is the new work: adaptive congruence recovery,
centered-root elimination, exceptional-case exclusion, positivity, and
canonical reconstruction. The second is the explicitly retained compiler
and calibration. No finite computation establishes that second equivalence
by searching for all possible roots.

\begin{longtable}{@{}p{0.27\textwidth}p{0.63\textwidth}@{}}
\toprule
Input or verification & Exact role\\
\midrule
\endhead
Source 44 compact proof & Whole-tuple native compiler adapter, mask bounds,
canonical padding and finite coefficient recipe. Retained verbatim.\\
Inherited ordinary proof and Pell audit & Native consistency kernel,
power/half-binomial recovery, canonical normalized witnesses, and named
K215/refutation calibration. Retained with source pins.\\
Sources 58 and 61 & Squared-modulus auxiliary reselection, positive gaps,
transport carrier, and signed index lemma. Retained and explicitly used.\\
New adaptive-carrier proof & Replaces $q^5,q^4$ by $q^2Y,qY$; handles the
factor three and proves the same existence equivalence.\\
New centered-gap proof & Merges $f,u$ into $n$; proves integrality,
positivity, the absence of the exceptional branch, and the converse.\\
Exact finite checkers & Polynomial identities, leading forms, degree
witnesses, gate liveness/counts, whole coefficient hashes and reconstruction.\\
Formalization status & No newly compiled end-to-end Lean/Rocq proof of
the changed literal polynomial is claimed.\\
\bottomrule
\end{longtable}

The primary new checks use only the Python standard library. They expand
the low-degree carrier residuals in all witnesses and symbolic ports, and
verify the large norm through exact identities in independent abstract
registers. This avoids an unnecessary full expansion of the degree-$110$
polynomial while still proving the relevant polynomial identity. The
nonzero coefficients in Section~\ref{sec:degree} certify degree after
specialization to the actual huge fixed numerals.

Separate small modular and bounded tests are diagnostics for transcription
and the factor-three argument. They are labeled as such. They are not
evidence of consistency, semantic completeness, or global minimality.
Independent checks in this archive were conducted by separate AI-assisted
workers in this research session; they are not external human peer review.

\subsection{Archive layout and reproducibility}

The main commands are
\begin{verbatim}
python3 verify.py
pdflatex -interaction=nonstopmode -halt-on-error report.tex
pdflatex -interaction=nonstopmode -halt-on-error report.tex
\end{verbatim}

\begin{itemize}
\item \texttt{report.tex}, \texttt{report.pdf}: this self-contained article.
\item \texttt{build.py}, \texttt{build\_domains.py}: regenerate the
twelve-variable primary, sum-of-squares, and domain companion circuits.
\item \texttt{positive13/}: generator, literal circuit, full proof, and
independent exact checker for the thirteen-variable reduction.
\item \texttt{circuits/}: authoritative division-free graphs for the
twelve-variable forms and standard domain companions.
\item \texttt{coefficients/}: four actual numeral payloads and the complete
unchanged compact44 binding.
\item \texttt{evaluate.py}: evaluate any delivered circuit at supplied
integers, optionally modulo a chosen integer; no zero search.
\item \texttt{checks/}, \texttt{VERIFICATION.json}: finite verification
scripts and receipts stating their exact scope.
\item \texttt{upstream/}: pinned proof dependencies, coefficient generator,
literal transition table, and license/attribution records.
\item \texttt{research/}: independent mathematical audit, rejected scale
shortcut, and the low-degree exploration with its corrected counts.
\end{itemize}

The new files can be placed as a contribution beside the existing sources.
The current smaller-variable and lower-operation alternatives should remain
available. This report establishes two new measured points and their proof
obligations; it does not identify a unique simplest equation across all
domains, degrees, circuit sizes, and coefficient heights.

\begin{thebibliography}{9}
\bibitem{source61}
V. Reshetnikov et al., ProveIt, source 61, \emph{Congruence carriers},
snapshot \texttt{a2cd79b9120}, 23 October 2026.
\href{https://github.com/VladimirReshetnikov/ProveIt/tree/a2cd79b9120fc75f48dcbef6e38d0f8b18d64d49/Computability/DiophantineConsistency/con-zfc-diophantine-forms/61-congruence-carriers}{Pinned source directory}.
The proof and baseline README are retained under \texttt{upstream/}.

\bibitem{source44}
ProveIt, source 44, \emph{Complement-packed native clauses: a smaller complete
coefficient tuple}, same snapshot.
\href{https://github.com/VladimirReshetnikov/ProveIt/tree/a2cd79b9120fc75f48dcbef6e38d0f8b18d64d49/Computability/DiophantineConsistency/con-zfc-diophantine-forms/44-b4-sparse/native}{Pinned native tuple and proof}.

\bibitem{source58}
ProveIt, source 58, \emph{Integral root}, same snapshot.
\href{https://github.com/VladimirReshetnikov/ProveIt/tree/a2cd79b9120fc75f48dcbef6e38d0f8b18d64d49/Computability/DiophantineConsistency/con-zfc-diophantine-forms/58-integral-root}{Pinned source directory}.
The report source is retained as \texttt{upstream/SOURCE58\_REPORT.tex}.

\bibitem{source62}
ProveIt, source 62, \emph{Sign and norm}, same snapshot.
\href{https://github.com/VladimirReshetnikov/ProveIt/tree/a2cd79b9120fc75f48dcbef6e38d0f8b18d64d49/Computability/DiophantineConsistency/con-zfc-diophantine-forms/62-sign-and-norm}{Pinned source directory}.

\bibitem{source64}
ProveIt, source 64, \emph{Structural reductions}, same snapshot.
\href{https://github.com/VladimirReshetnikov/ProveIt/tree/a2cd79b9120fc75f48dcbef6e38d0f8b18d64d49/Computability/DiophantineConsistency/con-zfc-diophantine-forms/64-structural-reductions}{Pinned source directory}.

\bibitem{source65}
ProveIt, source 65, \emph{Shared taps}, same snapshot.
\href{https://github.com/VladimirReshetnikov/ProveIt/tree/a2cd79b9120fc75f48dcbef6e38d0f8b18d64d49/Computability/DiophantineConsistency/con-zfc-diophantine-forms/65-shared-taps}{Pinned source directory}.
Its Pell component uses successive congruence recovery.

\bibitem{bayer}
J. Bayer, M. David, M. Ha{\ss}ler, Y. Matiyasevich, and D. Schleicher,
\emph{Diophantine equations over $\Z$: universal bounds and parallel
formalization}, arXiv:2506.20909, v2, 30 June 2025, Lemma 6.1.
\url{https://arxiv.org/html/2506.20909v2}.
The latest version observed during this work is v3, 22 October 2026;
the cited elementary domain lemma is pinned to v2.
\end{thebibliography}
\end{document}
